\originalchapter{两体弛豫、相混合与快速集体演化}

本章对应 8.902 第2讲末至第3讲. 要把恒星视作连续的无碰撞分布, 须比较轨道穿越时间与粒子离散引力引起的弛豫时间. “无碰撞”并非“无引力”: 恒星不断受共同引力势加速, 只是单独近距离遭遇改变轨道的速度通常较慢. 星系仍可在时变集体势中很快调整, 这种过程需要与两体随机扩散区分.

\originalsection{一次弱引力遭遇的冲量}
先取质量 $m$ 的场恒星近似静止, 另一恒星以相对速度 $v$ 沿直线 $x=vt$ 通过, 冲量参数为 $b$. 被测恒星的横向加速度大小
\[
a_\perp=\frac{Gmb}{(b^2+v^2t^2)^{3/2}}.
\]
积分奇对称的纵向分量为零, 横向分量为
\[
\delta v_\perp=\int_{-\infty}^{\infty}a_\perp\dd t
=\frac{Gm}{bv}\int_{-\infty}^{\infty}\frac{\dd u}{(1+u^2)^{3/2}}
=\frac{2Gm}{bv}.
\]
这个公式依赖 $\delta v\ll v$, 不能把它外推成任意大角散射. 两体相对轨道的精确散射角满足 $\tan(\Theta/2)=G(m_1+m_2)/(bv^2)$, 因而 $90^\circ$ 的尺度是 $b_{90}=G(m_1+m_2)/v^2$. 对等质量、维里尺度 $v^2\sim GN m/R$, 得 $b_{90}\sim2R/N$. 系数取决于这里的 $v$ 是相对速度还是单星速度, 后面估算只保留一致的数量级.

\begin{center}
\begin{tikzpicture}[>=Latex,scale=0.9]
\draw[->](-3,1.1)--(3,1.1) node[right]{$v$};
\fill(0,0)circle(2pt) node[below]{场恒星};
\draw[dashed](0,0)--(0,1.1) node[midway,right]{$b$};
\fill(-1.6,1.1)circle(2pt) node[above]{被测恒星};
\draw[->,structurecolor](-1.6,1.05)--(-1.3,0.65) node[left]{$a_\perp$};
\end{tikzpicture}
\end{center}
图用于直线弱偏转近似, 不画成精确双曲轨道. 在星系里 $b_{90}\ll R$, 强遭遇截面 $\pi b_{90}^2$ 很小. 均匀球的数密度 $n=3N/(4\pi R^3)$ 给一次穿越强遭遇的概率尺度 $n\pi b_{90}^2R\sim3/N$. 对 $N\sim10^{11}$ 极小; 对球状星团 $N\sim10^4$ 较大, 仍须累积许多次穿越才能估计整体演化.

\originalsection{许多弱遭遇产生对数积分}
在 $\dd t$ 内, 冲量参数在 $[b,b+\dd b]$ 的遭遇数近似 $\dd N_{
\rm enc}=2\pi b\dd b\,nv\dd t$. 若不同遭遇方向随机且相关可忽略, 平均冲量相消, 速度二次变化累积:
\[
\frac{\dd\langle(\Delta v)^2\rangle}{\dd t}
=\int_{b_{\min}}^{b_{\max}}2\pi bnv\left(\frac{2Gm}{bv}\right)^2\dd b
=\frac{8\pi G^2m^2n}{v}\ln\Lambda,
\quad\Lambda=\frac{b_{\max}}{b_{\min}}.
\]
每个对数区间的贡献相当, 因此大量远距离弱遭遇总效应可以超过罕见强遭遇. $b_{\min}$ 不能任意趋零: 应由大角偏转、恒星半径或其他小尺度物理截断; $b_{\max}$ 约为系统尺度或使均匀独立遭遇近似失效的尺度. $\ln\Lambda$ 被称作 Coulomb 对数, 名称来自相似积分结构, 此处相互作用仍是引力.

以速度改变达到原速度数量级定义 $t_{\rm rel}$, 得
\[
t_{\rm rel}\simeq\frac{v^3}{8\pi G^2m^2n\ln\Lambda}.
\]
用 $t_{\rm cross}=R/v$, $v^2\simeq GNm/R$ 及均匀球 $n$,
\[
\frac{t_{\rm rel}}{t_{\rm cross}}\simeq\frac{N}{6\ln\Lambda}
\sim\frac{N}{6\ln N}.
\]
系数6来自这一简化冲量与均匀球模型, 更完整的速度分布会改变系数. 物理结论由大的 $N/\ln N$ 决定. 星系 $R=10\,\mathrm{kpc}$、$v=200\,\mathrm{km\,s^{-1}}$ 的穿越时间约 $4.9\times10^7$ 年, 取 $N=10^{11}$ 给弛豫时间约 $3\times10^{16}$ 年, 远超课程宇宙年龄尺度. 无碰撞近似因而合理. 这个时间不是一条圆轨道周期, 后者另有 $2\pi$ 因子.

\begin{example}
简化球状星团取 $N=10^4$, $R=3\,\mathrm{pc}$, $v=5\,\mathrm{km\,s^{-1}}$, 有 $t_{\rm cross}\simeq0.59$ 百万年, $N/(6\ln N)\simeq181$, 故 $t_{\rm rel}$ 约一亿年. 这是一致尺度估算, 并未声称所有球状星团都有相同弛豫时间; 中心密度和质量谱会显著改变局部演化.
\end{example}

\originalsection{无碰撞分布仍可改变粗粒化外观}
若 $f(\boldsymbol x,\boldsymbol v,t)$ 是相空间数密度, 无碰撞动力学意味着沿轨道 $\dd f/\dd t=0$. 这并不表示固定位置的 $f$ 不变. 相空间团块可被轨道流拉伸和卷绕, 精细分布保持信息, 有限分辨的平均却趋于较平滑的外观. 将后者叫作粗粒化弛豫, 必须保留精细动力学的可逆性.

对可积定常势, 用作用量 $\boldsymbol J$ 与角变量 $\boldsymbol\theta$, 有 $\dot{\boldsymbol J}=0$, $\dot{\boldsymbol\theta}=\boldsymbol\Omega(\boldsymbol J)$. 不同 $J$ 的频率不同, 角向不均匀性逐渐失相. 一个 Fourier 分量的可观测平均包含
\[
\int\dd J\,a_k(J)\exp[-\ii k\Omega(J)t].
\]
若 $a_k$ 有有限宽度且 $\Omega'(J)\ne0$, 不同相位在积分中抵消, 混合时间量级 $t_{\rm mix}\sim[|k\Omega'|\Delta J]^{-1}$. 这是相混合. 若势严格谐振且所有轨道频率相同, 团块只共同转动, 不会凭这一机制变成均匀角分布. 因而“等一段时间必然相混合”需要频率展宽条件.

没有碰撞时精细 $f$ 的最大值也不会被动力学任意增加. 从光滑观测恢复一个粗粒化 $\bar f$ 常会丢失这些细丝, 但这属于分辨率与平均的选择, 不能当作已证明熵在每条精细轨道上严格增加.

\originalsection{时变引力势直接交换单星能量}
令单星的单位质量能量 $E=v^2/2+\Phi(\boldsymbol x,t)$, 运动方程 $\dot{\boldsymbol v}=-\nabla\Phi$. 用链式法则,
\[
\frac{\dd E}{\dd t}=\boldsymbol v\cdot\dot{\boldsymbol v}
+\frac{\partial\Phi}{\partial t}+\boldsymbol v\cdot\nabla\Phi
=\frac{\partial\Phi}{\partial t}.
\]
定常势保留每颗星的能量, 时变势却可交换它们的能量而无需近距离碰撞. 若势井在恒星通过时变深, 出井时可能带走较少能量; 变浅时可带走较多. 在自洽系统中恒星能量之和还需计入共享势能的配对计数, 不能简单把所有单位质量 $E$ 相加代替系统总机械能.

星系形成、塌缩和并合使集体势在动力学时间内大幅改变, 由此产生快速的暴烈弛豫（violent relaxation）. 可用
\[
t_{\rm vr}\sim\left\langle\left(\frac{\partial_t\Phi}{E}\right)^2\right\rangle^{-1/2}
\]
估计势变化引起的能量调整时间, 在非零典型 $E$ 的模型中它常与塌缩或自由落体时间同阶. 这不是普适精确公式, 也没有保证最终分布一定 Maxwell 化. 调整可能不完整, 还保留形成史与守恒量.

三种过程至此有清楚分工: 两体弛豫是离散遭遇的随机扩散; 相混合是在定常势中由频率差造成角相位失相; 暴烈弛豫依靠时变集体势改变单星能量. 后面以 Poisson 方程求共同势, 再以无碰撞 Boltzmann 方程和分布函数描述恒星, 正是建立在这些时间尺度比较上.
